\section{Lựa chọn mô hình triển khai Gemini}

\subsection{Vì sao gọi API, không tự host}

Tự host LLM đòi GPU, lượng hóa, giám sát VRAM --- ngoài phạm vi đồ án ngành một
sinh viên không có máy chủ GPU ổn định. Gemini Flash qua lớp OpenAI-compatible
cho phép đổi \texttt{base\_url} sang vLLM sau này mà không viết lại node
LangGraph. Đánh đổi: phụ thuộc mạng, hạn mức, chính sách nhà cung cấp, và không
nộp điểm thi mô hình mở.

\subsection{Vì sao Flash, không Pro}

Pro chậm và đắt hơn trên cùng số token filter+generate+HyDE. Flash đủ cho câu
tiếng Việt ngắn trên $k=8$ Điều. Nhiệt độ 0 giảm khác biệt giữa hai lần hỏi
cùng câu --- tiện demo hội đồng.

\subsection{Tách khóa nhúng và khóa sinh câu}

Hai biến môi trường cho phép sau này dùng nhà nhúng khác (rẻ, ổn định chiều)
và nhà sinh câu khác. Hiện tại có thể trùng một khóa AI Studio. Reindex chỉ cần
khóa nhúng; chat cần cả hai.

\section{Nhật ký, quan sát và gỡ lỗi}

\texttt{tool\_calls} trong state ghi tên nút, provider, số kết quả, lỗi. SSE đẩy
metadata rút gọn (top 5 retrieval) để UI hiện ``đang tìm'', không đẩy nguyên
nội dung 8 Điều (nặng). Debug chi tiết nằm server log, không trả client --- tránh
lộ prompt hệ thống.

Khi demo hỏng, thứ tự kiểm: container healthy không, Postgres có hàng luật
không, Chroma manifest khớp không, \texttt{.env} gốc có khóa không, CORS có
chặn dev không. Phụ lục A lặp lại thứ tự này cho người chấm.

\section{Sao lưu volume và máy trắng}

Ba volume cần giữ nếu muốn khôi phục dữ liệu demo: Postgres (tài khoản, hội
thoại, corpus chữ, lịch), Chroma (vector), uploads (file hợp đồng). Xóa volume
Postgres mà giữ Chroma làm lệch id Điều --- phải reindex. Sao lưu mức đồ án:
\texttt{docker compose down} không gắn \texttt{-v}; copy thư mục volume nếu cần
mang máy. Không cấu hình sao lưu tự động.

Máy trắng cài Docker, clone mã, điền \texttt{.env}, \texttt{compose up}, nạp
luật. Không đòi cài sẵn Python 3.10 hay Node trên host, trừ khi phát triển ngoài
Compose. WSL2 trên Windows chỉ là host; sản phẩm là Linux container.
